\documentclass[11pt, letterpaper]{article}
% ----- PAQUETES -----
\usepackage[utf8]{inputenc}
\usepackage[T1]{fontenc}
\usepackage[spanish]{babel}
\usepackage{graphicx}
\usepackage[margin=2.3cm]{geometry}
\usepackage{xcolor}
\usepackage{titlesec}
\usepackage{booktabs}
\usepackage{enumitem}
\usepackage{xurl}
\usepackage{hyperref}
\usepackage{fancyhdr}
\usepackage{amsmath}
\usepackage{amssymb}
\usepackage{array}
\usepackage{tikz}
\usetikzlibrary{arrows.meta, babel, calc, decorations.pathmorphing}

% ----- COLORES -----
\definecolor{UnalGreen}{RGB}{0,153,51}
\definecolor{DarkGray}{RGB}{60,60,60}
\definecolor{StripA}{RGB}{200,60,60}
\definecolor{StripB}{RGB}{60,90,180}

% ----- SECCIONES -----
\titleformat{\section}
{\normalfont\large\bfseries\color{UnalGreen}}
{\thesection}{1em}{}[{\titlerule[0.8pt]}]
\titleformat{\subsection}
{\normalfont\normalsize\bfseries\color{DarkGray}}
{\thesubsection}{1em}{}
\titlespacing*{\section}{0pt}{8pt}{4pt}
\titlespacing*{\subsection}{0pt}{6pt}{2pt}

% ----- ENLACES -----
\hypersetup{
	colorlinks=true,
	linkcolor=UnalGreen,
	urlcolor=UnalGreen
}

% ----- ENCABEZADO -----
\pagestyle{fancy}
\fancyhf{}
\lhead{\footnotesize Mecánica y Ondas para Biociencias (1000023) -- Laboratorio}
\rhead{\footnotesize Semestre 2026-2}
\cfoot{\thepage}

% ----- LISTAS COMPACTAS -----
\setlist[itemize]{leftmargin=*, itemsep=2pt, topsep=2pt, parsep=0pt}
\setlist[enumerate]{leftmargin=*, itemsep=3pt, topsep=3pt, parsep=0pt}

% ----- CAJA DE NOTA -----
\newcommand{\notabox}[1]{%
	\begin{center}
		\fcolorbox{UnalGreen}{UnalGreen!8}{%
			\parbox{0.94\textwidth}{\small #1}%
		}
	\end{center}
}

% ----- TECHO PARA LA FIGURA DE LOS RESORTES -----
\newcommand{\techo}[2]{%
	\fill[gray!40] (#1,5.2) rectangle (#2,5.45);
	\draw[very thick] (#1,5.2) -- (#2,5.2);%
}

\begin{document}
	
	% =========================================================
	% ENCABEZADO
	% =========================================================
	\begin{center}
		\vspace*{-0.8cm}
		\IfFileExists{imagenes/universidad-nacional-de-colombia-sede-bogota-seeklogo.png}{%
			\includegraphics[width=0.15\textwidth]
			{imagenes/universidad-nacional-de-colombia-sede-bogota-seeklogo.png}\\[0.1cm]
		}{}
		
		{\Large\bfseries Práctica 4: Ley de Hooke}\\[0.1cm]
		{\large\bfseries Rigidez de un resorte, en serie y en paralelo}\\[0.1cm]
		{\normalsize Mecánica y Ondas para Biociencias (1000023)}
	\end{center}
	
	\vspace{0.2cm}
	\noindent
	\begin{tabular}{@{}ll@{}}
		Modalidad: & Grupal (3 a 4 integrantes) \\
		Entregable: & Bitácora 4 \\
		Fecha de sesión: & 13 de octubre de 2026 \\
	\end{tabular}
	
	\notabox{\textbf{Antes de venir al laboratorio:} lean la descripción del
		experimento y respondan las \emph{preguntas previas}. El quiz individual del
		inicio de la sesión sale de esas preguntas.}
	
	% =========================================================
	\section{Objetivos}
	% =========================================================
	\begin{itemize}
		\item Caracterizar la rigidez $k$ de un resorte a partir de la relación
		entre fuerza y elongación.
		\item Explorar cómo se combinan las rigideces de dos resortes conectados
		en paralelo y en serie.
		\item Propagar la incertidumbre de una resta y representarla en una
		gráfica con barras de error.
	\end{itemize}
	
	% =========================================================
	\section{Materiales}
	% =========================================================
	\begin{itemize}
		\item Dos resortes de diferente rigidez.
		\item Juego de masas y portapesas.
		\item Soporte universal.
		\item Regla o cinta métrica.
		\item Una varilla o regla liviana (para el montaje en paralelo).
		\item Papel milimetrado.
	\end{itemize}
	
	% =========================================================
	\section{Descripción del experimento}
	% =========================================================
	
	Un resorte se estira cuando se cuelga una masa de él. Se mide la longitud
	inicial $L_0$ del resorte sin carga y, para cada masa $m$ colgada, su
	longitud $L_i$; la \textbf{elongación} es $\Delta L = L_i - L_0$.
	
	\vspace{0.1cm}
	\noindent
	El experimento se hace con tres configuraciones: cada resorte
	\textbf{individual}, los dos resortes \textbf{en paralelo} (lado a lado,
	sosteniendo juntos la masa) y los dos \textbf{en serie} (uno a continuación
	del otro).
	
	\vspace{0.3cm}
	\begin{center}
		\begin{tikzpicture}[scale=0.9, every node/.style={font=\small}]
			\tikzset{resorte/.style={decorate,
					decoration={zigzag, segment length=3mm, amplitude=1.4mm,
						pre=lineto, pre length=3mm, post=lineto, post length=3mm},
					thick, UnalGreen!70!black}}
			\tikzset{masa/.style={thick, fill=gray!25}}
			
			% ---------- (a) individual ----------
			\techo{0}{2}
			\draw[resorte] (1,5.2) -- (1,3.2);
			\draw[masa] (0.55,2.5) rectangle (1.45,3.2);
			\node at (1,2.85) {$m$};
			\node[right] at (1.15,4.2) {$k$};
			\node at (1,1.7) {(a) Individual};
			
			% ---------- (b) paralelo ----------
			\techo{3.6}{7.0}
			\draw[resorte] (4.4,5.2) -- (4.4,3.8);
			\draw[resorte] (6.2,5.2) -- (6.2,3.8);
			\draw[very thick] (4.1,3.8) -- (6.5,3.8);
			\draw[thick] (5.3,3.8) -- (5.3,3.2);
			\draw[masa] (4.85,2.5) rectangle (5.75,3.2);
			\node at (5.3,2.85) {$m$};
			\node[left] at (4.3,4.5) {$k_1$};
			\node[right] at (6.3,4.5) {$k_2$};
			\node at (5.3,1.7) {(b) Paralelo};
			
			% ---------- (c) serie ----------
			\techo{8.2}{10.2}
			\draw[resorte] (9.2,5.2) -- (9.2,4.0);
			\fill (9.2,4.0) circle (1.6pt);
			\draw[resorte] (9.2,4.0) -- (9.2,2.9);
			\draw[masa] (8.75,2.2) rectangle (9.65,2.9);
			\node at (9.2,2.55) {$m$};
			\node[right] at (9.35,4.6) {$k_1$};
			\node[right] at (9.35,3.45) {$k_2$};
			\node at (9.2,1.7) {(c) Serie};
		\end{tikzpicture}
	\end{center}
	
	% =========================================================
	\section{Preguntas previas}\label{sec:previas}
	% =========================================================
	
	Respóndanlas \textbf{antes} de la sesión. No se califican, pero el quiz de
	entrada sale de aquí.
	
	\begin{enumerate}
		\item \emph{Predicción} (respondan antes de hacer cálculos): si se cuelgan
		masas cada vez mayores de un resorte, ¿qué forma tiene la gráfica de la
		elongación $\Delta L$ contra la masa? Y si se cuelga la misma masa de dos
		resortes en paralelo, ¿se estira más o menos que con uno solo? ¿Y en
		serie?
		
		\item Enuncien la ley de Hooke. ¿Qué mide la constante $k$ y en qué
		unidades? ¿Qué resorte es más rígido: el de $k$ grande o el de $k$
		pequeño?
		
		\item \emph{Paralelo.} Los dos resortes se estiran lo \textbf{mismo},
		$\Delta L$, y la fuerza total es la suma de las dos fuerzas. Demuestren
		que $k_{\text{par}} = k_1 + k_2$.
		
		\item \emph{Serie.} La \textbf{misma} fuerza $F$ actúa sobre los dos
		resortes y la elongación total es la suma de las dos elongaciones.
		Demuestren que
		\[
		\frac{1}{k_{\text{ser}}} = \frac{1}{k_1} + \frac{1}{k_2}
		\]
	\end{enumerate}
	
	% =========================================================
	\section{Teoría}
	% =========================================================
	
	Dentro de su región elástica, la fuerza que hay que aplicar a un resorte es
	proporcional a su elongación:
	\begin{equation}
		F = k\,\Delta L
		\label{eq:hooke}
	\end{equation}
	
	\noindent
	Cuando el resorte cuelga en equilibrio con una masa $m$, la fuerza aplicada
	es el peso de la masa, $F = m\,g$. Graficando $F$ (eje vertical) contra
	$\Delta L$ (eje horizontal) se obtiene una recta cuya \textbf{pendiente es
		$k$}.
	
	\notabox{Si grafican $\Delta L$ contra $m$ (como en otras guías), la
		pendiente es $g/k$ y \textbf{no} $k$. Por eso aquí graficamos $F$ contra
		$\Delta L$: así la pendiente es directamente la constante del resorte.}
	
	\subsection*{Combinación de resortes}
	
	De las preguntas previas 3 y 4 se obtiene:
	\begin{equation}
		k_{\text{par,teo}} = k_1 + k_2
		\label{eq:paralelo}
	\end{equation}
	\begin{equation}
		k_{\text{ser,teo}} = \frac{1}{\dfrac{1}{k_1} + \dfrac{1}{k_2}}
		= \frac{k_1\,k_2}{k_1 + k_2}
		\label{eq:serie}
	\end{equation}
	
	\noindent
	En paralelo la rigidez \textbf{aumenta} (el sistema es más difícil de
	estirar); en serie \textbf{disminuye} (queda más blando que el más blando de
	los dos).
	
	\subsection*{Incertidumbres}
	
	La elongación es una \textbf{resta} de dos longitudes medidas con la misma
	regla. Con la regla de la suma y resta del Taller~2:
	\begin{equation}
		\Delta(\Delta L) = \sqrt{(\Delta L_i)^2 + (\Delta L_0)^2}
		= \sqrt{2}\,\Delta_{\text{reg}}
		\label{eq:dDL}
	\end{equation}
	
	\noindent
	donde $\Delta_{\text{reg}}$ es la resolución de la regla. Para la fuerza,
	$\Delta F = g\,\Delta m$, con $\Delta m$ la resolución de las masas.
	
	% =========================================================
	\section{Procedimiento}
	% =========================================================
	
	\begin{enumerate}
		\item Cuelguen un resorte del soporte y midan su longitud sin carga, $L_0$.
		\item Cuelguen masas crecientes (al menos \textbf{cinco}) y registren la
		longitud $L_i$ con cada una. Esperen a que el resorte se quede quieto
		antes de leer.
		\item Calculen $\Delta L = L_i - L_0$ en cada caso.
		\item Repitan con el \textbf{otro resorte}.
		\item Repitan con los dos resortes en \textbf{paralelo}: cuelguen ambos
		del soporte, unan sus extremos inferiores con la varilla y cuelguen la
		masa del \textbf{centro} de la varilla, para que los dos se estiren por
		igual. Midan $L_0$ de este montaje sin carga.
		\item Repitan con los dos resortes en \textbf{serie}: el segundo resorte
		cuelga del extremo del primero, y la masa del extremo del segundo.
		La longitud $L$ es la del conjunto completo.
	\end{enumerate}
	
	\notabox{\textbf{No sobrecarguen los resortes.} Usen solo masas con las que
		el resorte regrese a su longitud original al retirarlas. Si queda
		estirado de forma permanente, ya salió de la región elástica y la ley de
		Hooke deja de aplicar.}
	
	\vspace{0.2cm}
	\begin{center}
		\small
		$\Delta_{\text{reg}} = $ \rule{1.5cm}{0.4pt} cm
		\qquad
		$\Delta m = $ \rule{1.5cm}{0.4pt} g
		\qquad
		$L_0 = $ \rule{1.5cm}{0.4pt} cm
		\quad (una tabla por configuración)
		
		\vspace{0.3cm}
		\begin{tabular}{|c|c|c|c|c|}
			\hline
			$m$ (g) & $F = mg$ (N) & $L_i$ (cm) & $\Delta L$ (cm)
			& $\Delta(\Delta L)$ (cm) \\
			\hline
			\rule{0pt}{3ex} & & & & \\ \hline
			\rule{0pt}{3ex} & & & & \\ \hline
			\rule{0pt}{3ex} & & & & \\ \hline
			\rule{0pt}{3ex} & & & & \\ \hline
			\rule{0pt}{3ex} & & & & \\ \hline
			\rule{0pt}{3ex} & & & & \\ \hline
		\end{tabular}
	\end{center}
	
	% =========================================================
	\section{Análisis de resultados}
	% =========================================================
	
	\subsection{Fuerza y elongación}
	
	Para cada configuración, calculen $F = m\,g$ \textbf{con $m$ en kg}, usando
	$g = 9.78$ m/s$^2$, y $\Delta(\Delta L)$ con la ecuación (\ref{eq:dDL}).
	Verifiquen si $\Delta F$ es pequeño frente a la incertidumbre de $\Delta L$:
	si lo es, basta con dibujar barras de error \emph{horizontales}.
	
	\subsection{Las gráficas}
	
	Para cada una de las cuatro configuraciones (resorte 1, resorte 2, paralelo y
	serie), grafiquen $F$ (eje vertical, en N) contra $\Delta L$ (eje
	horizontal, \textbf{en metros}), con barras de error. Pueden hacerlo en una
	sola gráfica con colores distintos para compararlas. Tracen a ojo la recta
	que mejor represente los puntos de cada una.
	
	\subsection{Las constantes}
	
	Obtengan la pendiente de cada recta, escogiendo dos puntos sobre la recta
	(no sobre los datos) bien separados, como en el Taller~2:
	\[
	k = \frac{F_2 - F_1}{\Delta L_2 - \Delta L_1}
	\]
	
	\noindent
	Reporten $k_1$, $k_2$, $k_{\text{par}}$ y $k_{\text{ser}}$ en N/m.
	
	\subsection{Comparación con la teoría}
	
	Calculen los valores teóricos con las ecuaciones (\ref{eq:paralelo}) y
	(\ref{eq:serie}), usando sus $k_1$ y $k_2$, y comparen:
	\[
	\text{disc.\ paralelo} = \frac{|k_{\text{par}} - (k_1+k_2)|}{k_1+k_2}\times 100\%
	\qquad
	\text{disc.\ serie} = \frac{|k_{\text{ser}} - k_{\text{ser,teo}}|}{k_{\text{ser,teo}}}\times 100\%
	\]
	
	\begin{center}
		\small
		\begin{tabular}{|l|c|c|c|}
			\hline
			Configuración & $k$ medida (N/m) & $k$ teórica (N/m) & Discrepancia \\
			\hline
			\rule{0pt}{3ex} Resorte 1 & & --- & --- \\ \hline
			\rule{0pt}{3ex} Resorte 2 & & --- & --- \\ \hline
			\rule{0pt}{3ex} Paralelo & & & \\ \hline
			\rule{0pt}{3ex} Serie & & & \\ \hline
		\end{tabular}
	\end{center}
	
	% =========================================================
	\section{Conclusiones}
	% =========================================================
	
	Además de sus conclusiones propias, discutan:
	
	\begin{itemize}
		\item ¿Los resortes cumplen la ley de Hooke? ¿Sus puntos se alinean, o
		la gráfica se curva a cargas mayores?
		\item ¿La recta de cada configuración pasa por el origen? ¿Debería
		hacerlo? ¿Qué podría explicar un intercepto distinto de cero (por
		ejemplo, el peso del propio resorte)?
		\item ¿Qué configuración resultó más rígida y cuál más blanda? ¿Coincide
		con su predicción de la pregunta previa 1?
		\item ¿Las expresiones teóricas para serie y paralelo describen bien sus
		datos? ¿Las discrepancias son comparables con la dispersión de sus
		puntos?
		\item ¿Qué aspectos del montaje (lectura de $L$, balanceo de la masa,
		varilla desnivelada en paralelo) pueden explicar las diferencias?
	\end{itemize}
	
	% =========================================================
	\section{Referencias}
	% =========================================================
	
	\begin{itemize}
		\item L. Agulles Perdós, \emph{Guías de Laboratorio: Mecánica y Ondas
			para Biociencias}, Departamento de Física, Universidad Nacional de
		Colombia, 2026. Laboratorio 10.
		\item C. Perilla, \emph{Guía de trabajo en el laboratorio: Fuerza de un
			resorte}, Universidad Nacional de Colombia, 2023.
	\end{itemize}
	
	\vfill
	\noindent
	{\footnotesize \textit{Auxiliar docente: Miguel Andrés Perdomo Gutiérrez
			--- mperdomog@unal.edu.co}}
	
\end{document}